% Options for packages loaded elsewhere
\PassOptionsToPackage{unicode}{hyperref}
\PassOptionsToPackage{hyphens}{url}
\documentclass[
  ignorenonframetext,
]{beamer}
\newif\ifbibliography
\usepackage{pgfpages}
\setbeamertemplate{caption}[numbered]
\setbeamertemplate{caption label separator}{: }
\setbeamercolor{caption name}{fg=normal text.fg}
\beamertemplatenavigationsymbolsempty
% remove section numbering
\setbeamertemplate{part page}{
  \centering
  \begin{beamercolorbox}[sep=16pt,center]{part title}
    \usebeamerfont{part title}\insertpart\par
  \end{beamercolorbox}
}
\setbeamertemplate{section page}{
  \centering
  \begin{beamercolorbox}[sep=12pt,center]{section title}
    \usebeamerfont{section title}\insertsection\par
  \end{beamercolorbox}
}
\setbeamertemplate{subsection page}{
  \centering
  \begin{beamercolorbox}[sep=8pt,center]{subsection title}
    \usebeamerfont{subsection title}\insertsubsection\par
  \end{beamercolorbox}
}
% Prevent slide breaks in the middle of a paragraph
\widowpenalties 1 10000
\raggedbottom
\AtBeginPart{
  \frame{\partpage}
}
\AtBeginSection{
  \ifbibliography
  \else
    \frame{\sectionpage}
  \fi
}
\AtBeginSubsection{
  \frame{\subsectionpage}
}
\usepackage{iftex}
\ifPDFTeX
  \usepackage[T1]{fontenc}
  \usepackage[utf8]{inputenc}
  \usepackage{textcomp} % provide euro and other symbols
\else % if luatex or xetex
  \usepackage{unicode-math} % this also loads fontspec
  \defaultfontfeatures{Scale=MatchLowercase}
  \defaultfontfeatures[\rmfamily]{Ligatures=TeX,Scale=1}
\fi
\usepackage{lmodern}
\ifPDFTeX\else
  % xetex/luatex font selection
\fi
% Use upquote if available, for straight quotes in verbatim environments
\IfFileExists{upquote.sty}{\usepackage{upquote}}{}
\IfFileExists{microtype.sty}{% use microtype if available
  \usepackage[]{microtype}
  \UseMicrotypeSet[protrusion]{basicmath} % disable protrusion for tt fonts
}{}
\makeatletter
\@ifundefined{KOMAClassName}{% if non-KOMA class
  \IfFileExists{parskip.sty}{%
    \usepackage{parskip}
  }{% else
    \setlength{\parindent}{0pt}
    \setlength{\parskip}{6pt plus 2pt minus 1pt}}
}{% if KOMA class
  \KOMAoptions{parskip=half}}
\makeatother
\setlength{\emergencystretch}{3em} % prevent overfull lines
\providecommand{\tightlist}{%
  \setlength{\itemsep}{0pt}\setlength{\parskip}{0pt}}
% Slide layout defaults for warning-clean scientific decks.
\usepackage{etoolbox}
\IfFileExists{xurl.sty}{\usepackage{xurl}}{}
\IfFileExists{seqsplit.sty}{\usepackage{seqsplit}}{\newcommand{\seqsplit}[1]{#1}}
\protected\def\breakseq#1{\seqsplit{#1}}
\protected\def\breaktt#1{\begingroup\ttfamily\seqsplit{#1}\endgroup}
\setlength{\emergencystretch}{6em}
\tolerance=5000
\hbadness=10000
\hfuzz=1pt
\setlength{\tabcolsep}{2pt}
\AtBeginEnvironment{longtable}{\tiny\renewcommand{\arraystretch}{0.86}\setlength{\tabcolsep}{1pt}}
\AtBeginEnvironment{tabular}{\tiny\renewcommand{\arraystretch}{0.86}\setlength{\tabcolsep}{1pt}}

% Natbib and cross-reference fallbacks — slides don't load natbib
% or cleveref, but combined-PDF manuscript prose may emit these
% commands. The fallback renders citations as a bracketed key list
% and cross-references as detokenized labels so slides don't fail on
% undefined control sequences or raw underscores. \providecommand is
% a no-op if packages load later.
\providecommand{\citep}[1]{[#1]}
\providecommand{\citet}[1]{#1}
\providecommand{\citealp}[1]{#1}
\providecommand{\citeauthor}[1]{#1}
\providecommand{\citeyear}[1]{#1}
\providecommand{\cref}[1]{\texttt{\detokenize{#1}}}
\providecommand{\Cref}[1]{\texttt{\detokenize{#1}}}

\newtheorem{proposition}{Proposition}
\newtheorem{hypothesis}{Hypothesis}
\usepackage{bookmark}
\IfFileExists{xurl.sty}{\usepackage{xurl}}{} % add URL line breaks if available
\urlstyle{same}
% fallback for those not using the hyperref driver hyperxmp:
\makeatletter
\@ifundefined{xmpquote}{\newcommand{\xmpquote}[1]{#1}}{}
\makeatother
\hypersetup{
  hidelinks,
  pdfcreator={LaTeX via pandoc}}

\author{\texorpdfstring{}{}}
\date{}

\begin{document}

\section{Bibliometric and Temporal
Analysis}\label{bibliometric-and-temporal-analysis}

\begin{frame}[allowframebreaks]{Bibliometric and Temporal Analysis}
Descriptive statistics summarize the corpus along every available axis:
counts by year, venue, and author; citation-count distributions; and
author productivity. Temporal analysis fits the publication time series,
reporting a compound annual growth rate of 3.45\% across 2000--2026 (a
span of 26 years), with a mean year-over-year growth rate of 6.3\% and a
doubling time of 11.3 years. The peak publication year is 2025 with 112
records.
\end{frame}

\begin{frame}[fragile,allowframebreaks]{Growth Metrics}
\protect\phantomsection\label{growth-metrics}
The compound annual growth rate (CAGR) is computed as:

\[
\text{CAGR} = \left(\frac{N_{\text{end}}}{N_{\text{start}}}\right)^{1/(\text{year span})} - 1
\]

where \(N_{\text{start}}\) is the publication count in the first year
(2000) and \(N_{\text{end}}\) is the count in the last year (2026). The
mean year-over-year growth rate \(\bar{g}\) is the arithmetic mean of
annual ratios. The doubling time is
\(t_d = \ln(2) / \ln(1 + \text{CAGR})\). These metrics are stored in
\breaktt{temporal\_analysis.json} and injected into the manuscript at
render time.
\end{frame}

\begin{frame}[fragile,allowframebreaks]{Subfield Classification}
\protect\phantomsection\label{subfield-classification}
Subfield classification assigns each record to one of 6 configurable
buckets (Clinical Sleep, Cognition, Pharmacology, Psychiatry, Safety,
and Neuroscience) by priority-aware keyword matching; the taxonomy is
defined entirely in configuration
(\breaktt{project\_config.subfield\_keywords}). The largest bucket is
\textbf{Clinical Sleep} at 64.3\% of the classified corpus. A
per-subfield temporal breakdown (\breaktt{subfield\_timeline.json})
tracks how each sub-area has grown over time, enabling identification of
emerging or declining research threads.
\end{frame}

\begin{frame}[allowframebreaks]{Topic Modeling}
\protect\phantomsection\label{topic-modeling}
A TF-IDF term-weighting of titles and abstracts \breakseq{[@salton1988term]}
feeds non-negative matrix factorization (NMF) \breakseq{[@lee1999learning]},
implemented with scikit-learn \breakseq{[@pedregosa2011scikit]}. NMF decomposes
the document-term matrix \(\mathbf{V} \approx \mathbf{W} \mathbf{H}\),
where \(\mathbf{W}\) is the document-topic matrix and \(\mathbf{H}\) is
the topic-term matrix. The factorization extracts 5 latent topics that
cross-cut the keyword taxonomy. The random seed is fixed at 42 for
reproducibility. The reporting follows established systematic-review
practice \breakseq{[@page2021prisma]}, with every figure and statistic
traceable to a committed artifact.
\end{frame}

\end{document}
